14
9 *Continuous Mappings (General Theory)
squares (with sides parallel to the axes) are the products of open intervals, which are
open sets in R.
It should be noted that the sets $G_1 \times G_2$, where $G_1 \in \tau_1$ and $G_2 \in \tau_2$, constitute
only a base for the topology, not all the open sets in the direct product of topological
spaces.

9.2.4 Problems and Exercises
1. Verify that if $(X, d)$ is a metric space, then $\left(X, \frac{d}{1+d}\right)$ is also a metric space, and
the metrics $d$ and $\frac{d}{1+d}$ induce the same topology on $X$. (See also Problem 1 of the
preceding section.)
2. a) In the set $\mathbb{N}$ of natural numbers we define a neighborhood of the number
$n \in \mathbb{N}$ to be an arithmetic progression with difference $d$ relatively prime to $n$. Is the
resulting topological space Hausdorff?
b) What is the topology of $\mathbb{N}$, regarded as a subset of the set $\mathbb{R}$ of real numbers
with the standard topology?
c) Describe all open subsets of $\mathbb{R}$.
3. If two topologies $\tau_1$ and $\tau_2$ are defined on the same set, we say that $\tau_2$ is stronger
than $\tau_1$ if $\tau_1 \subset \tau_2$, that is $\tau_2$ contains all the sets in $\tau_1$ and some additional open sets
not in $\tau_1$.
a) Are the two topologies on $\mathbb{N}$ considered in the preceding problem compara-
ble?
b) If we introduce a metric on the set $C[0, 1]$ of continuous real-valued functions
defined on the closed interval $[0, 1]$ first by relation (9.6) of Sect. 9.1, and then by
relation (9.7) of the same section, two topologies generally arise on $C[a, b]$. Are
they comparable?
4. a) Prove in detail that the space of germs of continuous functions defined in
Example 4 is not Hausdorff.
b) Explain why this topological space is not metrizable.
c) What is the weight of this space?
5. a) State the axioms for a topological space in the language of closed sets.
b) Verify that the closure of the closure of a set equals the closure of the set.
c) Verify that the boundary of any set is a closed set.
d) Show that if $F$ is closed and $G$ is open in $(X, \tau)$, then the set $G \setminus F$ is open in
$(X, \tau)$.
e) If $(Y, \tau_Y)$ is a subspace of the topological space $(X, \tau)$, and the set $E$ is such
that $E \subset Y \subset X$ and $E \in \tau_X$, then $E \in \tau_Y$.
6. A topological space $(X, \tau)$ in which every point is a closed set is called a topo-
logical space in the strong sense or a $t_1$-space. Verify the following statements.
a) Every Hausdorff space is a $t_1$-space (partly for this reason, Hausdorff spaces
are sometimes called $t_2$-spaces).